\chapter[Límits i colímits]{Límits i colímits}

Aquest test recull els conceptes essencials del capítol. Cada pregunta té una única resposta correcta.

\begin{Form}
\iniciTest{lim}{b,c,b,c,d,b,c,b,d,c,b,c,b,c,b}

\begin{enumerate}

\pregunta Què és un \textbf{diagrama} de forma $\cat{I}$ en una categoria $\cat{C}$?
\begin{enumerate}
  \opcio{a}{Un objecte de $\cat{C}$.}
  \opcio{b}{Un functor $D\colon\cat{I}\to\cat{C}$, on $\cat{I}$ és una categoria índex petita.}
  \opcio{c}{Una transformació natural.}
  \opcio{d}{Un morfisme de $\cat{I}$.}
\end{enumerate}
\feedback

\pregunta Un \textbf{con} sobre $D$ amb vèrtex $A$ és el mateix que:
\begin{enumerate}
  \opcio{a}{un objecte terminal.}
  \opcio{b}{un límit.}
  \opcio{c}{una transformació natural $\Delta_A\Rightarrow D$ des del functor constant.}
  \opcio{d}{un functor $\cat{C}\to\cat{I}$.}
\end{enumerate}
\feedback

\pregunta Un \textbf{límit} de $D$ és:
\begin{enumerate}
  \opcio{a}{qualsevol con sobre $D$.}
  \opcio{b}{un objecte terminal de la categoria de cons $\mathrm{Con}(D)$.}
  \opcio{c}{un objecte inicial de $\cat{C}$.}
  \opcio{d}{el producte de tots els objectes.}
\end{enumerate}
\feedback

\pregunta Quina categoria índex dona lloc al \textbf{producte} de dos objectes?
\begin{enumerate}
  \opcio{a}{La categoria buida.}
  \opcio{b}{Dues fletxes paral·leles.}
  \opcio{c}{La categoria discreta amb dos objectes.}
  \opcio{d}{La forma de racó.}
\end{enumerate}
\feedback

\pregunta A $\cat{Set}$, l'\textbf{equalitzador} de $f,g\colon A\to B$ és:
\begin{enumerate}
  \opcio{a}{el producte cartesià $A\times B$.}
  \opcio{b}{la unió disjunta.}
  \opcio{c}{el quocient de $B$.}
  \opcio{d}{el subconjunt $\{a\in A\mid f(a)=g(a)\}$.}
\end{enumerate}
\feedback

\pregunta El \textbf{pullback} de $f\colon A\to C$ i $g\colon B\to C$ a $\cat{Set}$ és:
\begin{enumerate}
  \opcio{a}{$A\times B$ sencer.}
  \opcio{b}{$\{(a,b)\in A\times B\mid f(a)=g(b)\}$.}
  \opcio{c}{la unió $A\cup B$.}
  \opcio{d}{el quocient $A/B$.}
\end{enumerate}
\feedback

\pregunta L'\textbf{objecte terminal} és el límit de quin diagrama?
\begin{enumerate}
  \opcio{a}{El diagrama de dues fletxes paral·leles.}
  \opcio{b}{El diagrama d'un sol objecte.}
  \opcio{c}{El diagrama buit.}
  \opcio{d}{El diagrama de racó.}
\end{enumerate}
\feedback

\pregunta Un \textbf{colímit} d'un diagrama $D$ és:
\begin{enumerate}
  \opcio{a}{un altre nom per al límit.}
  \opcio{b}{un objecte inicial a la categoria de cocons (un límit a la categoria oposada).}
  \opcio{c}{sempre l'objecte terminal.}
  \opcio{d}{un con.}
\end{enumerate}
\feedback

\pregunta El \textbf{coproducte} a $\cat{Set}$ és:
\begin{enumerate}
  \opcio{a}{el producte cartesià.}
  \opcio{b}{la intersecció.}
  \opcio{c}{el conjunt buit.}
  \opcio{d}{la unió disjunta.}
\end{enumerate}
\feedback

\pregunta Quina és la caracterització dels \textbf{límits finits}?
\begin{enumerate}
  \opcio{a}{N'hi ha prou amb tenir coproductes.}
  \opcio{b}{Cal tenir tots els límits de forma petita.}
  \opcio{c}{Equivalen a tenir objecte terminal, productes binaris i equalitzadors (o terminal i pullbacks).}
  \opcio{d}{No es poden construir a partir de casos més simples.}
\end{enumerate}
\feedback

\pregunta Quan \textbf{preserva} un functor $F$ el límit d'un diagrama $D$?
\begin{enumerate}
  \opcio{a}{Sempre, per definició de functor.}
  \opcio{b}{Quan $(F(L),F\lambda)$ és un límit de $F\comp D$.}
  \opcio{c}{Quan $F$ és constant.}
  \opcio{d}{Quan $F$ no té adjunts.}
\end{enumerate}
\feedback

\pregunta El functor oblit $U\colon\cat{Grp}\to\cat{Set}$:
\begin{enumerate}
  \opcio{a}{no preserva ni límits ni colímits.}
  \opcio{b}{preserva colímits però no límits.}
  \opcio{c}{preserva límits però no colímits.}
  \opcio{d}{preserva tot.}
\end{enumerate}
\feedback

\pregunta El functor representable $\Hom(A,-)$ sempre:
\begin{enumerate}
  \opcio{a}{no és un functor.}
  \opcio{b}{preserva els límits que existeixen.}
  \opcio{c}{envia límits a colímits.}
  \opcio{d}{és constant.}
\end{enumerate}
\feedback

\pregunta Per què límits i colímits són tan útils?
\begin{enumerate}
  \opcio{a}{Perquè només existeixen a $\cat{Set}$.}
  \opcio{b}{Perquè no tenen propietats universals.}
  \opcio{c}{Perquè unifiquen productes, coproductes, equalitzadors, pullbacks i objectes terminal/inicial com a objectes universals d'un diagrama.}
  \opcio{d}{Perquè eviten els functors.}
\end{enumerate}
\feedback

\pregunta En el diagrama següent, quina propietat universal fa de $P$ el \textbf{pullback} de $f$ i $g$?
\[
\begin{tikzpicture}[baseline=(m.center)]
  \matrix (m) [matrix of math nodes, row sep=2.6em, column sep=2.6em]
    { P & A \\ B & C \\ };
  \draw[-{Stealth[scale=1.1]}] (m-1-1) -- node[above] {$p_1$} (m-1-2);
  \draw[-{Stealth[scale=1.1]}] (m-1-1) -- node[left] {$p_2$} (m-2-1);
  \draw[-{Stealth[scale=1.1]}] (m-1-2) -- node[right] {$f$} (m-2-2);
  \draw[-{Stealth[scale=1.1]}] (m-2-1) -- node[below] {$g$} (m-2-2);
\end{tikzpicture}
\]
\begin{enumerate}
  \opcio{a}{Que $f\comp p_1=g\comp p_2$ i res més.}
  \opcio{b}{Que per a tot objecte $X$ amb fletxes $x_1\colon X\to A$, $x_2\colon X\to B$ tals que $f\comp x_1=g\comp x_2$, hi ha un únic $u\colon X\to P$ amb $p_1\comp u=x_1$ i $p_2\comp u=x_2$.}
  \opcio{c}{Que $P=A\times B$ sempre.}
  \opcio{d}{Que $f$ i $g$ són isomorfismes.}
\end{enumerate}
\feedback

\end{enumerate}
\clauestatica
\finalitzaTest
\end{Form}
